\documentclass[11pt,a4paper]{article}

\usepackage[T1]{fontenc}
\usepackage[utf8]{inputenc}
\usepackage{lmodern}
\usepackage{microtype}
\usepackage[a4paper,margin=1in]{geometry}
\usepackage{amsmath,amssymb,bm}
\usepackage[hidelinks]{hyperref}

\newcommand{\dd}{\,\mathrm{d}}
\newcommand{\R}{\mathbb{R}}
\newcommand{\bx}{\bm{x}}
\newcommand{\bc}{\bm{c}}
\newcommand{\bn}{\bm{n}}
\newcommand{\bu}{\bm{u}}
\newcommand{\bU}{\bm{U}}
\newcommand{\bOmega}{\bm{\Omega}}
\newcommand{\bomega}{\bm{\omega}}
\newcommand{\bI}{\bm{I}}
\newcommand{\bL}{\bm{L}}
\newcommand{\bLambda}{\bm{\Lambda}}
\newcommand{\bp}{\bm{p}}
\newcommand{\bh}{\bm{h}}
\newcommand{\bP}{\bm{P}}
\newcommand{\bH}{\bm{H}}
\newcommand{\btau}{\bm{\tau}}
\newcommand{\bzero}{\bm{0}}
\newcommand{\quat}[1]{\mathbf{#1}}
\DeclareMathOperator{\vecop}{vec}

\title{Impulse Time Integration Subsection for POISE}
\author{}
\date{\today}

\begin{document}
\maketitle
\setcounter{section}{3}
\setcounter{subsection}{2}

\subsection{Impulse time integration}
\label{sec:impulse-time-integration}

The boundary-element solve described above gives the instantaneous velocity
potential $\phi$ on the body surfaces.  For the production potential-flow
sweeps in this work we do not evaluate hydrodynamic force from a finite
difference approximation to $\partial\phi/\partial t$.  Instead, $\phi$ is used
to compute Lamb impulse state functions, and the rigid-body update is written
in impulse-difference form.  This avoids storing a temporal history of the
potential and gives excellent conservation of total linear momentum in the
multi-body calculations.

Consider $N_b$ rigid bodies.  Body $b$ has centroid $\bc_b(t)\in\R^3$,
translational velocity $\bU_b(t)\in\R^3$, lab-frame angular velocity
$\bOmega_b(t)\in\R^3$, mass $m_b$, body-frame inertia tensor $\bI_b$, and
surface $S_b$.  The normal $\bn$ follows the convention used in the boundary
integral equation, and the no-penetration condition supplies the Neumann data
\begin{equation}
  \frac{\partial \phi}{\partial n}
  =
  \bn\cdot
  \left[
    \bU_b+\bOmega_b\times(\bx-\bc_b)
  \right],
  \qquad \bx\in S_b .
  \label{eq:impulse-neumann}
\end{equation}
Once the BEM system has been solved, the linear and angular Lamb impulses of
body $b$ are evaluated by the surface integrals
\begin{align}
  \bL_b
  &=
  \rho_f\int_{S_b}\phi\,\bn\,\dd S,
  \label{eq:linear-lamb-impulse}\\
  \bLambda_b
  &=
  \rho_f\int_{S_b}
    \phi\,(\bx-\bc_b)\times\bn\,\dd S .
  \label{eq:angular-lamb-impulse}
\end{align}
Here $\rho_f$ is the fluid density.  The angular impulse
$\bLambda_b$ is taken about the translating centroid $\bc_b$, not about a fixed
origin.

The impulse variables entering the rigid-body dynamics are
\begin{align}
  \bp_b
  &=
  m_b\bU_b-\bL_b,
  \label{eq:body-linear-impulse}\\
  \bh_b
  &=
  \bc_b\times\bp_b+\bh_{s,b}-\bLambda_b,
  \label{eq:body-angular-impulse}
\end{align}
where $\bp_b$ is the contribution of body $b$ to the total linear impulse and
$\bh_b$ is its contribution to the total angular impulse.  The solid angular
momentum is
\begin{equation}
  \bh_{s,b}
  =
  \vecop\!\left[
    \quat{q}_b
    \left(0,\bI_b\bomega_b\right)
    \quat{q}_b^{-1}
  \right],
  \qquad
  (0,\bomega_b)=
  \quat{q}_b^{-1}(0,\bOmega_b)\quat{q}_b ,
  \label{eq:solid-angular-impulse}
\end{equation}
with $\quat{q}_b$ the unit quaternion mapping body-frame vectors into the
laboratory frame.  The total impulses are therefore
\begin{equation}
  \bP=\sum_{b=1}^{N_b}\bp_b,\qquad
  \bH=\sum_{b=1}^{N_b}\bh_b .
  \label{eq:total-impulses}
\end{equation}
In the absence of external forcing, the continuous inviscid problem conserves
$\bP$ and $\bH$.

Let $\Delta t=t^{n+1}-t^n$.  At the beginning of a time step the BEM solve gives
$\bL_b^n$ and $\bLambda_b^n$.  The endpoint velocity and torque are found by a
fixed-point iteration because the endpoint impulses depend on the unknown
endpoint configuration.  Given a trial endpoint velocity $\bU_b^{n+1}$ and
endpoint angular state, the centroid is advanced by the trapezoidal rule
\begin{equation}
  \bc_b^{n+1}
  =
  \bc_b^n+
  \frac{\Delta t}{2}
  \left(\bU_b^n+\bU_b^{n+1}\right).
  \label{eq:impulse-position-update}
\end{equation}
The BEM is then solved again at the trial endpoint to obtain
$\bL_b^{n+1}$ and $\bLambda_b^{n+1}$.  With no optional energy-gradient or
pressure correction, the translational residual used by the code is
\begin{equation}
  \bm{r}_{U,b}
  =
  m_b\left(\bU_b^{n+1}-\bU_b^n\right)
  -
  \left(\bL_b^{n+1}-\bL_b^n\right).
  \label{eq:linear-impulse-residual}
\end{equation}
Thus convergence of the iteration enforces
\begin{equation}
  m_b\bU_b^{n+1}-\bL_b^{n+1}
  =
  m_b\bU_b^n-\bL_b^n,
  \label{eq:per-body-linear-impulse}
\end{equation}
which is the discrete per-body linear-impulse balance used in the fast impulse
scheme.  The Newton-like correction to $\bU_b^{n+1}$ is preconditioned in the
body frame by $m_b\mathbf{1}+\mathbf{M}_{a,b}$, where $\mathbf{M}_{a,b}$ is the
diagonal single-body added-mass estimate used only as a preconditioner.

The angular update must account for the fact that
\eqref{eq:angular-lamb-impulse} is measured about the moving centroid.  Define
\begin{equation}
  \bp_{b}^{m}
  =
  \frac{1}{2}
  \left[
    \left(m_b\bU_b^n-\bL_b^n\right)
    +
    \left(m_b\bU_b^{n+1}-\bL_b^{n+1}\right)
  \right],
  \qquad
  \bU_b^m=\frac{1}{2}\left(\bU_b^n+\bU_b^{n+1}\right).
  \label{eq:midpoint-impulse}
\end{equation}
The step-averaged torque applied to the solid angular momentum is
\begin{equation}
  \btau_b^m
  =
  \frac{\bLambda_b^{n+1}-\bLambda_b^n}{\Delta t}
  -
  \bU_b^m\times\bp_b^m .
  \label{eq:impulse-torque}
\end{equation}
The second term is the transport correction for angular momentum about a
translating point.  Omitting it gives the wrong angular-impulse balance even for
a single body.  In the converged linear solve, $\bp_b^m=\bp_b^n=\bp_b^{n+1}$,
and \eqref{eq:impulse-torque} is the discrete counterpart of
\[
  \frac{\dd\bh_{s,b}}{\dd t}
  =
  \frac{\dd\bLambda_b}{\dd t}
  -
  \bU_b\times\bp_b .
\]

The orientation and angular velocity are advanced with an implicit-midpoint
Lie-group update.  Let $\bomega_b^n$ be the body-frame angular velocity at the
start of the step.  The midpoint body-frame angular velocity solves
\begin{equation}
  \bomega_b^m
  =
  \bomega_b^n
  +
  \frac{\Delta t}{2}
  \bI_b^{-1}
  \left[
    \btau_b^m(\quat{q}_b^m)
    -
    \bomega_b^m\times\left(\bI_b\bomega_b^m\right)
  \right],
  \label{eq:implicit-midpoint-omega}
\end{equation}
where $\btau_b^m(\quat{q}_b^m)$ denotes the lab-frame torque
\eqref{eq:impulse-torque} rotated into the midpoint body frame.  The midpoint
orientation is
\begin{equation}
  \quat{q}_b^m
  =
  \exp\left[
    \frac{\Delta t}{2}(0,\bOmega_b^m)
  \right]\quat{q}_b^n,
  \label{eq:midpoint-orientation}
\end{equation}
and the accepted endpoint state is reconstructed as
\begin{align}
  \quat{q}_b^{n+1}
  &=
  \exp\left[
    \Delta t(0,\bOmega_b^m)
  \right]\quat{q}_b^n,
  \label{eq:endpoint-orientation}\\
  \bomega_b^{n+1}
  &=
  2\bomega_b^m-\bomega_b^n .
  \label{eq:endpoint-angular-velocity}
\end{align}
The exponential map in \eqref{eq:midpoint-orientation} and
\eqref{eq:endpoint-orientation} is evaluated as
\[
  \exp(0,\bm{a})
  =
  \left(\cos\|\bm{a}\|,\,
  \frac{\sin\|\bm{a}\|}{\|\bm{a}\|}\bm{a}\right),
\]
with the usual limiting value when $\|\bm{a}\|=0$.

The fixed-point iteration is applied to the stacked endpoint translational
velocities and step-averaged torques, with Anderson acceleration used to reduce
the iteration count.  At convergence, the method has advanced the bodies using
only BEM state functions evaluated at the start and end of the step.  In the
production runs reported here no algebraic energy projection is applied:
kinetic energy is monitored as a diagnostic of the impulse discretisation,
while linear impulse is conserved to near roundoff except in runs that terminate
early because the state becomes non-finite.

\end{document}
